%%%%%%%%%%%%%%%%%%%%%%%%%%%%%%%%%%%%%%%%%%%%%%%%%%%%%%%%%%%%%%%%%%%%%%
% Lớp 10 | Hệ thức lượng trong tam giác | Định lý sin và định lý côsin
% Dạng 0H4?2-3: Biến đổi, rút gọn biểu thức
% Số câu: 10  (sinh tự động bởi tach_cau.py; từ nay thầy sửa trực tiếp file này)
%%%%%%%%%%%%%%%%%%%%%%%%%%%%%%%%%%%%%%%%%%%%%%%%%%%%%%%%%%%%%%%%%%%%%%

%%%%%%%%%%%%%%%% TRẮC NGHIỆM 4 PHƯƠNG ÁN %%%%%%%%%%%%%%%%

% nguồn: MaHoaNH_TracNghiem_Lop10
\begin{ex}%[0H4B2-3]%[ma:5ecb895a]
	Trong mặt  phẳng $O x y$, cho đường tròn $(\mathscr{C}):(x-7)^2+(y+3)^2=100$. Phương trình tiếp tuyến của $(\mathscr{C})$ tại điểm $A(1; 5)$ có dạng $a x+b y+17=0$. Giá trị của $a^3-b^2$ là
	\choice 
		{$ -55 $} 
		 {$ 25 $}
		{$ 43  $}
		{\True$  11 $}
	\loigiai{$(\mathscr{C}):(x-7)^2+(y+3)^2=100$ có tâm $ I (7;-3)  $.\\
		Tiếp tuyến tại điểm  $A(1; 5)$ nhận $ \overrightarrow{IA}=(-6;8) $ là véc-tơ pháp tuyến.\\
	Do  đó phương trình tiếp tuyến là $ -6(x-1)+8(y-5)=0 $, tức là $ 3x-4y+17=0 $.}
\end{ex}

% nguồn: MaHoaNH_TracNghiem_Lop10
\begin{ex}%[0H4B2-3]%[ma:0ed14500]
	Trong mặt phẳng $Oxy$, cho  đường thẳng $(\Delta)\colon 4x-3y-3=0$. Đường tròn tâm $I(1;-3)$ nào sau đây tiếp xúc với $(\Delta)$?
	 \choice
		{$(x+1)^2+(y-3)^2=4$} 
		{\True  $(x-1)^2+(y+3)^2=4$}
		 {$(x+1)^2+(y-3)^2=2$}
		{$(x-1)^2+(y+3)^2=2$} 
	\loigiai{ 
		Ta có đường tròn tâm $I(1;-3)$ tiếp xúc với $(\Delta)$ nên bán kính đường  tròn là \[R=\mathrm{d}(I,\Delta)=\dfrac{\left|4-3(-3)-3\right|}{\sqrt{4^2+(-3)^2}}=2.\]
		Phương  trình đường tròn $(C)\colon (x-1)^2+(y+3)^2=4$.
	} 
\end{ex}

% nguồn: MaHoaNH_TracNghiem_Lop10
\begin{ex}%[0H4B2-3]%[ma:9d6c2790]
	Trong mặt phẳng tọa độ $Oxy$, cho đường tròn $(C)\colon x^2+y^2+2x-6y+5=0$. Gọi $\Delta$ là tiếp tuyến của $(C)$ tại điểm $A(0;1)$. Phương  trình tổng quát của $\Delta$ là
	\choice 
		{$x-2y+1=0$} 
		{$x+2y-2=0$} 
		{$x+2y+2=0$} 
		{\True $x-2y+2=0$} 
	 \loigiai{
		Đường  tròn $(C)$ có tâm $I(-1;3)$.\\
		Phương trình tiếp tuyến $\Delta$ của $(C)$ tại $A(0;1)$ nhận véc-tơ $\overrightarrow{IA}=(1;-2)$ làm véc-tơ pháp  tuyến có dạng
		 $$\Delta\colon 1(x-0)-2(y-1)=0\Leftrightarrow x-2y+2=0.$$
	 }
\end{ex}

% nguồn: MaHoaNH_TracNghiem_Lop10
\begin{ex}%[0H4B2-3]%[ma:eba190e5]
	Cho đường tròn  $(C)\colon x^2+y^2-3 x-y=0$. Phương trình tiếp tuyến của $(C)$ tại $M(1;-1)$ là:
	 \choice
		{$x-3  y-2=0$}
		{\True $x+3 y+2=0$} 
		 {$x-3 y+2=0$}
		{$x+3  y-2=0$}
	 \loigiai{
		Đường tròn $ (C) $ có tâm $ I\left(\dfrac{3}{2};\dfrac{1}{2}\right) $.\\ 
		 Ta có $ \overrightarrow{MI}=\left(\dfrac{1}{2};\dfrac{3}{2}\right) =\dfrac{1}{2}(1;3)$.\\
		Do đó phương trình tiếp tuyến của đường tròn $  (C) $ tại $ M $ là $ 1(x-1)+3(y+1)=0\Leftrightarrow x+3y+2=0 $.
	} 
\end{ex}

% nguồn: MaHoaNH_TracNghiem_Lop10
\begin{ex}%[0H4B2-3]%[ma:a559dc89]
	Trong mặt phẳng tọa độ $Oxy$, cho đường tròn $(C)$ có  phương trình $(x-1)^2+(y+2)^2=8$. Phương trình tiếp tuyến $d$ của đường tròn $(C)$ tại điểm $A(3;-4)$ là
	\choice 
		{$d\colon  x+y+1=0$}
		{$d\colon  x-2y-11=0$}
		{\True  $d\colon x-y-7=0$}
		{$d\colon  x-y+7=0$}
	\loigiai{ 
		Đường tròn $(C)$ có tâm $I(1;-2)$.\\ 
		Ta  có $ \overrightarrow{IA}=(2;-2) =2\cdot(1;-1)$.\\
		 Tiếp tuyến $d$ của đường tròn $(C)$ tại điểm $A(3;-4)$ có một véc-tơ pháp tuyến là $ \overrightarrow{n}=(1;-1) $
		Phương trình tiếp tuyến $d$  của đường tròn $(C)$ tại điểm $A(3;-4)$ là \[d\colon (x-3)-(y+4)=0\Leftrightarrow x-y-7=0.\]
	} 
\end{ex}

% nguồn: MaHoaNH_TracNghiem_Lop10
\begin{ex}%[0H4B2-3]%[ma:c564a9b5]
	Đường tròn nào sau đây tiếp  xúc với trục $Ox$?
	 \choice
		{$x^2+y^2-5=0$} 
		 {$x^2+y^2-10x-2y+1=0$}
		{\True $x^2+y^2+6x+5y+9=0$} 
		{$x^2+y^2-10x=0$} 
	\loigiai{ 
	 Phương trình trục $Ox$ là $y=0$. \\
	Xét đường tròn $(C)\colon x^2+y^2+6x+5y+9=0$ có  tâm $I\left(-3;-\dfrac{5}{2}\right)$ và bán kính $R=\dfrac{5}{2}$. \\
	Ta  có $\mathrm{d}(I,Ox)=\left|-\dfrac{5}{2}\right|=\dfrac{5}{2} \Rightarrow \mathrm{d}(I,Ox)=R$. \\
	Vậy đường tròn $(C)\colon x^2+y^2+6x+5y+9=0$ tiếp  xúc với trục $Ox$.
	 }
\end{ex}

% nguồn: MaHoaNH_TracNghiem_Lop10
\begin{ex}%[0H4K2-3]%[ma:e6884b19]
	Trong mặt phẳng $Oxy$, cho đường tròn $(C)\colon (x-1)^2+(y-4)^2=4$. Phương trình tiếp tuyến với đường tròn $(C)$  song song với đường thẳng $\Delta\colon 4x-3y-2=0$ là
	 \choice
		 {$4x-3y-18=0$}
		{$4x-3y+18=0;4x-3y-2=0$} 
		{$4x-3y-18=0;4x-3y+2=0$} 
		 {\True $4x-3y+18=0$}
	\loigiai{ 
		Đường tròn $(C)\colon (x-1)^2+(y-4)^2=4$ có tâm $I(1;4)$ và  bán kính $r=2$.\\
		Tiếp tuyến $d$ của $(C)$ song song với đường thẳng $\Delta\colon  4x-3y-2=0$ có phương trình dạng
		$$d\colon  4x-3y+c=0,\; c\neq -2.$$
		Ngoài  ra $\mathrm{d}(I,d)=r \Leftrightarrow \dfrac{\big|4\cdot 1-3\cdot 4+c\big|}{\sqrt{4^2+(-3)^2}}=2
		 \Leftrightarrow \big|c-8\big|=10 \Leftrightarrow \hoac{&c=-2\;\;\text{(loại)}\\&c=18\;\;\text{(nhận).}}$\\
		Vậy $d\colon 4x-3y+18=0$  là tiếp tuyến cần tìm.
	} 
\end{ex}


%%%%%%%%%%%%%%%% ĐÚNG/SAI %%%%%%%%%%%%%%%%

% nguồn: MaHoaNH_DungSai_Lop10
\begin{ex}%[0H4H2-3]%[ma:063d6746]
	Cho tam giác $ABC$ không cân  thoả mãn $b^2-c^2=2a(b-c)$. Mỗi kết quả dưới đây đúng hay sai?
	\choiceTF 
		 {\True $b+c=2a$}
		{$b=8$ với $c=5$, $a=6$} 
		{\True $\sin B+\sin C=2\sin  A$}
		 {\True $b(\cos C-2)=c(\cos B-2)$}
	 \loigiai{
		\begin{itemchoice}
			\itemch Ta có $b^2-c^2=2a(b-c)\Leftrightarrow (b-c)(b+c)=2a(b-c)\Leftrightarrow b+c=2a$. (do tam giác không cân  nên $b\ne c$)
			\itemch Từ a) ta có $b+5=2\cdot 6\Leftrightarrow b=7$. 
			\itemch Từ a) và theo định lý sinta có $2R\sin B+2R\sin C=2\cdot 2R\sin A\Leftrightarrow \sin B+\sin  C=2\sin A$.
			\itemch Ta  có $b^2=a^2+c^2-2ac\cos B$ và $c^2=a^2+b^2-2ab\cos C$ nên
			\begin{eqnarray*}
				&  & b^2-c^2=c^2-b^2-2ac\cos B+2ab\cos C\\
				&\Leftrightarrow& 2(b^2-c^2)=2a(b\cos  C-c\cos B)\\
				&\Leftrightarrow& b^2-c^2=a(b\cdot  \cos C-c\cdot \cos B).
			\end{eqnarray*}
			Suy ra $a(b\cdot \cos C-c\cdot \cos B)=2a(b-c)\Leftrightarrow b\cdot \cos C-c\cdot \cos B=2b-2c\Leftrightarrow b(\cos  C-2)=c(\cos B-2)$.
		\end{itemchoice}
	 }
\end{ex}

% nguồn: MaHoaNH_DungSai_Lop10
\begin{ex}%[0H4H2-3]%[ma:bd195fef]
	Cho tam giác $ABC$ bất kỳ có $BC=a$, $AC=b$, $AB=c$. Mỗi kết quả dưới  đây đúng hay sai?
	\choiceTF 
		 {$b^2=a^2+c^2-2ac\cos A$}
		{\True $b^2-c^2=a(b\cdot \cos C-c\cdot \cos B)$} 
		{\True  $(b+c)\sin A=a(\sin B+\sin C)$}
		{\True $h_a=2R\cdot \sin B \cdot \sin C$} 
	 \loigiai{
		\begin{itemchoice}
			\itemch Theo đinh lý hàm số côsinthì  $b^2=a^2+c^2-2ac\cos B$.
			\itemch Ta có $b^2=a^2+c^2-2ac\cos  B$ và $c^2=a^2+b^2-2ab\cos C$ nên
			\begin{eqnarray*}
				& & b^2-c^2=c^2-b^2-2ac\cos  B+2ab\cos C\\
				&\Leftrightarrow&  2(b^2-c^2)=2a(b\cos C-c\cos B)\\
				&\Leftrightarrow& b^2-c^2=a(b\cdot  \cos C-c\cdot \cos B).
			\end{eqnarray*}
			\itemch Áp dụng hệ quả của  định lý sinta có
			\begin{eqnarray*}
				(b+c)\sin  A&=& (2R\sin B+2R\sin C)\cdot \dfrac{a}{2R}\\
				&=& (\sin B+\sin  C)\cdot \dfrac{2Ra}{2R}\\
				&=& a(\sin B+\sin  C).
			\end{eqnarray*}
			 \itemch Áp dụng hệ quả của định lý sinta có
			\[2R\cdot \sin B \cdot \sin C=2R\cdot \dfrac{b}{2R}\cdot  \dfrac{c}{2R}=\dfrac{bc}{2R}=\dfrac{abc}{4R}\cdot \dfrac{2}{a}=\dfrac{2S}{a}=h.\]
		\end{itemchoice}
	} 
\end{ex}


%%%%%%%%%%%%%%%% TỰ LUẬN %%%%%%%%%%%%%%%%

% nguồn: MaHoaNH_TuLuan_Lop10
\begin{bt}%[0H4H2-3]%[ma:b8673785]
\end{bt}
